\section{Conclusion}
This report presented GestureTalk, a real-time sign language recognition, text conversion and speech generation system built as a Project-II prototype. The application reads one hand from a laptop webcam, classifies a custom vocabulary of more than five hundred gestures, shows the recognised English word on the screen, and speaks that word through the Windows speech engine without using the internet.

The engineering work of this stage sits in the complete local pipeline rather than in a published accuracy number. MediaPipe Tasks supplies current landmark tracking. Geometry rules and named poses convert landmarks into words. Hysteresis, majority vote and a short confirmation window keep the overlay steady. A background speech thread lets a new pose interrupt the previous sentence. Synthetic tests cover named poses, handedness correction and noisy landmarks without a camera.

The system is a working demonstration, not a full interpreter of American Sign Language or Indian Sign Language. That limit is part of the conclusion, not an afterthought. Within the chosen scope, the project shows that a student implementation can go from camera to spoken word in real time, with a large everyday vocabulary, on an ordinary Windows laptop.

\section{Future Scope}
The next natural step is to collect landmark sequences from several people and train a small multilayer perceptron, keeping a person-wise split so that evaluation does not leak the same signer into both sides. A CustomTkinter shell can then wrap the OpenCV window, a history of recognised words, and simple controls for volume and handedness. ONNX export would allow a fair FPS and latency table on CPU. Two-handed poses, motion signs and a small phrase buffer are later extensions. Those items belong to later work. They are not claimed as results of the present prototype.
